\subsection{DET-14: actual period frames and exact continuation conditions}
The period-map route does supply positive finite operators independently
of zeros. This section constructs them from the actual source data
and isolates what the displayed finite-family realization can and
cannot provide for an infinite determinant.

\paragraph{Exact source maps, normalization and existence.}
In \path{Development/AbelianVarietyTheory/Definitions/Jets/SumIoPrAlt.lean},
\texttt{IsBeL}, lines 135--139, the normalized factorization is
\[
 \operatorname{be}(x,v)\eta=\langle l(v),\eta\rangle l(x),
 \qquad x\in R.\mathrm{sm},\quad v\in R.\mathrm{TK},
\]
with the paper's conjugate-first inner product.
\path{Development/AbelianVarietyTheory/Definitions/Jets/SideFam.lean},
lines 36--40, defines a common-target system:
for each integral shift \(\nu\) with \(n_1(\nu)\le q_0\),
the map \(l_\nu:(R_\nu).H\to W_p\) satisfies \texttt{IsBeL},
the number \(d_\nu=\mathrm{dg}(\nu)\) is positive, and it equals
the specified formal degree at every nonzero vector.
The shift has \(n=2r+k+1\) integer coordinates; this cutoff is finite,
since \(|\nu_i|\le q_0\).

This is supported by an actual existence export, not just a
definition of a desired factorization:
\path{CompanionTheory/AutomorphicTheory/PieriRules/ThetaBesselIsometries/ExistsFamilyFunctionalsStatement.lean},
lines 17--18, supplies left and right \texttt{SideFam} systems under
its complete \texttt{HypsG}, \texttt{Num}, \texttt{IsFamG},
parameter and measure hypotheses. It assumes finite-dimensional
\(W_p,W_q\). The underlying complete proof in
\path{CompanionTheory/AutomorphicTheory/PieriRules/ThetaBesselIsometries/ExistsFamilyFunctionals.lean},
lines 21--50, was read: it constructs normalized maps on the
bounded shift set, chooses zero outside it, and sets
\(d_\nu=c\,\mathrm{fdegHC}(\lambda+\nu)\), with \(c>0\) and
positivity on that set. We retain those representation-theoretic
hypotheses and dependencies; this source reading is not an independent
re-proof of all the imported analytic theorems.

\paragraph{DET-14(a): positive frames and their ordinary determinants.}
Assume such a \texttt{SideFam} system. Choose a finite set \(I\) of
admissible shifts, vectors
\(v_\nu\in(R_\nu).\mathrm{sm}\cap(R_\nu).\mathrm{TK}\), and
real \(t_\nu\ge0\). Define
\[
 a_\nu=\sqrt{t_\nu d_\nu}\,l_\nu(v_\nu),\qquad
 A_I=\sum_{\nu\in I}t_\nu d_\nu
                \operatorname{be}_\nu(v_\nu,v_\nu).
\]
Then \(A_I\) is positive and finite rank on the common target \(W_p\),
with
\[
 A_I=\sum_{\nu\in I}a_\nu a_\nu^*,\qquad
 \Tr A_I=\sum_{\nu\in I}t_\nu d_\nu\|l_\nu(v_\nu)\|^2.
\]
For the positive Gram matrix
\(M_I=(\langle a_\nu,a_\mu\rangle)_{\nu,\mu\in I}\),
\[
 \boxed{\quad
 \det_{W_p}(I-wA_I)=\det_{\mathbb C^I}(I-wM_I)
 =\sum_{J\subset I}(-w)^{|J|}
       \left\|\bigwedge_{\nu\in J}a_\nu\right\|^2.\quad}
\]
In particular, the alternating Taylor coefficients are nonnegative.
These are explicit frame determinants constructed from the source
maps; they have not thereby been identified with the historical
determinants of the Remark.
\begin{proof}
Both vector-domain requirements are imposed, so the source
factorization with \(x=v=v_\nu\) gives
\[
 t_\nu d_\nu\operatorname{be}_\nu(v_\nu,v_\nu)
 =a_\nu a_\nu^*,\qquad
 (aa^*)\eta=\langle a,\eta\rangle a.
\]
Each summand is bounded, selfadjoint and of rank at most one.
Moreover
\[
 \langle\eta,A_I\eta\rangle
 =\sum_{\nu\in I}|\langle a_\nu,\eta\rangle|^2\ge0.
\]
If \(a_\nu\ne0\), its rank-one summand has sole nonzero eigenvalue
\(\|a_\nu\|^2\); if \(a_\nu=0\), it contributes zero. Finite
additivity of the trace proves the stated trace identity.

Let \(S:\mathbb C^I\to W_p\) have columns \(a_\nu\).
Then \(A_I=SS^*\) and \(M_I=S^*S\).
For each nonzero eigenvalue \(\lambda\) of \(S^*S\), \(S\) maps its
eigenspace injectively onto the \(\lambda\)-eigenspace of \(SS^*\),
with inverse \(\lambda^{-1}S^*\). The converse follows in the
same way. These are positive selfadjoint finite-rank operators,
so their nonzero eigenvalues agree with multiplicities and there
are no nontrivial Jordan blocks. Extra zero eigenvalues contribute
only factors \(1\), proving the first determinant equality.

In the column expansion of \(\det(I-wM_I)\), choose the
\(-wM_I\) column exactly at positions \(J\).
The identity columns outside \(J\) force the corresponding rows
in the permutation expansion; the result is
\((-w)^{|J|}\det((M_I)_J)\), with no extra sign.
This principal minor is the Gram determinant of the vectors
\(a_\nu\), \(\nu\in J\). If they are dependent, both the determinant
and their wedge vanish. Otherwise Gram--Schmidt gives their
Gram determinant as the product of the squared lengths of successive
orthogonal residuals, which is precisely the squared Hilbert
exterior norm. For \(J=\varnothing\), both quantities are \(1\).
Summing over \(J\) proves the boxed expansion.
\end{proof}
The normalization \(\sqrt{d_\nu}\,l_\nu\) is also explicitly used in
\path{Development/AbelianVarietyTheory/Proofs/FlagMeasures/BesselExteriorPowerConstantBoundLargeGap.lean},
lines 80--85. The additional weights \(t_\nu\) and vectors \(v_\nu\)
are choices in our frame construction, not a source-supplied Xi
spectral choice. Zero weights and zero period vectors are permitted.

\paragraph{DET-14(b): necessary and sufficient trace-class condition.}
Suppose, after specified compatible isometric embeddings where needed,
the vectors above form a countable family \(a_1,a_2,\ldots\) in a
single Hilbert space \(\mathcal H\). Set
\(A_N=\sum_{j\le N}a_ja_j^*\). Then
\[
 \sum_j\|a_j\|^2<\infty
 \quad\Longleftrightarrow\quad
 \sup_N\Tr A_N<\infty
 \quad\Longleftrightarrow\quad
 A_N\ \hbox{converges in trace norm}.
\]
The limit \(A\) is positive and
\[
 \Tr A=\sum_j\|a_j\|^2,\qquad
 \|A-A_N\|_1=\sum_{j>N}\|a_j\|^2.
\]
\begin{proof}
The trace norm of \(a_ja_j^*\) is its nonnegative sole eigenvalue
\(\|a_j\|^2\). Thus summability makes the series absolutely
convergent in the Banach space \(S_1\), with sum \(A\).
Trace-norm convergence implies operator-norm convergence.
Limits of nonnegative quadratic forms are nonnegative, so \(A\)
and each tail \(A-A_N\) are positive.
The trace is continuous in trace norm, since
\(|\Tr T|\le\|T\|_1\). Therefore their traces are the corresponding
scalar sums. For a positive trace-class operator the trace norm
equals the trace, proving the exact tail formula.
The finite identity \(\Tr A_N=\sum_{j\le N}\|a_j\|^2\)
shows equivalence of summability and bounded partial traces.
Conversely, trace-norm convergence makes these continuous traces
bounded, proving the remaining implication.
\end{proof}
For these actual period vectors the condition is exactly
\[
 \sum_j t_jd_j\|l_j(v_j)\|^2<\infty,
\]
after the common-space compatibility has been supplied.
The positive rank-one construction makes this a trace-ideal
criterion rather than an inference from summable eigenvalues of a
nonnormal operator. By \fid{DET-05}, these conditions imply local
uniform convergence of the ordinary Fredholm determinants.
An identification with \(G\) still requires a separately proved
spectral or coefficient identity independent of selecting the
vectors or weights from Xi's zeros.

\paragraph{DET-14(c): the fixed target and fixed parameters cannot supply that limit.}
There are two precise restrictions on this displayed realization.
First, if the target \(W_p\) in the existence export is fixed of
dimension \(d\), then \(\operatorname{rank}A_I\le d\) and all its
determinants have degree at most \(d\), however many terms are
included. Any trace-norm limit acting on that target has the
same degree bound. More generally, any locally uniform limit on
\(\mathbb C\) of polynomials of degrees at most \(d\) has degree
at most \(d\): Cauchy's integral formula for derivatives passes
the zero \((d+1)\)-st derivatives to the limit. Thus such a limit
cannot be the infinite-zero Xi determinant \(G\).
This is a fixed-target restriction; a varying-target or
infinite-dimensional period realization is a different route.

Second, the exact \texttt{Num} hypothesis in
\path{Development/AbelianVarietyTheory/Definitions/Jets/SumIoPrAlt.lean},
lines 155--166, includes at line 165
\[
 100(q_0+1)\le ga,\qquad 100(q_0+1)\le h.
\]
Consequently
\[
 q_0\le\frac{\min(ga,h)}{100}-1.
\]
For fixed finite real \(ga,h\), the displayed existence theorem
cannot be invoked along \(q_0\to\infty\).
At least those data must vary, with any further changes necessary
to retain the remaining hypotheses. Compatibility of the resulting
spaces, representations and maps must then be established.
This is a direct implication of the written numerical hypotheses,
not an absence claim based on source searching.

\paragraph{The integrated normalization and the signed period identity.}
\path{CompanionTheory/AutomorphicTheory/PieriRules/ThetaBesselIsometries/FamilySetupBesselGrowthStatement.lean},
lines 16--17, asserts an \(L^2\) bound and, for a nonzero Bessel map,
the normalized isometry under its \texttt{HypsG} and
\texttt{SetupG} hypotheses:
\[
 d_\nu\int
 \|l_\nu(\sigma_\nu(g)v)\|^2\,d\nu_q(g)=\|v\|^2
\]
for the indicated smooth vectors. This provides useful common
\(L^2\) data when the parameters and measures are fixed; it is
a different realization from the finite target \(W_p\) just considered.
It does not by itself bound the sum of point values
\(\sum_jt_jd_j\|l_j(v_j)\|^2\).
To see why an integrated norm alone does not imply an evaluation
bound, take a continuous compactly supported scalar \(\phi\) with
\(\|\phi\|_2=1\), \(\phi(0)\ne0\), and let
\(f_n(x)=\sqrt n\,\phi(nx)\).
Substitution gives \(\|f_n\|_2=1\), whereas
\(|f_n(0)|=\sqrt n|\phi(0)|\to\infty\).
This example diagnoses the missing estimate; it does not assert
that those spike functions are the source Bessel coefficients.

The actual signed exterior-period decomposition also retains its
own coefficients. In
\path{Development/AbelianVarietyTheory/Proofs/FlagMeasures/BlockProjectionBesselStatement.lean},
lines 9--10, and its application in
\path{Development/AbelianVarietyTheory/Proofs/FlagMeasures/FirstNormAddSubsetIndicatorLe.lean},
lines 269--289, they are
\(\epsilon_I\kappa_{1,I}\overline{\kappa_{0,I}}d_I\).
Even identifying the two side maps to make the kappa product a
squared modulus does not remove the possible sign
\(\epsilon_I=-1\). Positive formal degrees alone therefore
do not identify that signed expansion with the Gram minors proved
above. Such an identification requires another specific identity.

The concrete upstream task is now to specify the varying source
data and compatible positive realization, prove the corresponding
trace bound, and identify its determinant with \(G\) independently
of the zero-based construction. These results do not provide the
separate historical correspondence to the Remark's \(W,S\) and
determinant family, and do not exclude another source construction.
